\section{관련 연구}
\label{sec:related}

\subsection{정간보의 구조}

정간보는 세로로 이어지는 각(刻)이 오른쪽에서 왼쪽으로 진행하고, 각은
정간(井間)이라 부르는 정사각형 칸의 열로 이루어진다. 정간 하나가 한 박이며,
여러 정간을 묶는 굵은 가로선인 대강(大綱)이 장단의 호흡 단위를 나타낸다
(예: 12박 = 3·3·3·3). 정간 안에는 12율(황종--응종)의 율명을 한자 또는
한글로 적고, 옥타브는 접두어(청·중청·배·하배) 또는 변형 한자(황종 黃의
청성 潢처럼 삼수변·인변을 더한 글자)로 구분한다. 한 정간 안에 음이 여럿이면 칸을 세로로 나눠 박을 등분하고(분박),
같은 행 안에 가로로 이어 쓰면 그 분박을 다시 나눈다(붙임). 즉 시가 정보가
전적으로 \emph{위치}에 인코딩된다는 점이 오선보와의 근본적 차이다
~\cite{kwon2000notation,kim2025omr}. 여기에 농음·퇴성·추성 등 연주법을
지시하는 시김새 기호가 음표 곁 또는 독립 칸에 부가되고, 장구 장단(구음),
가사가 별도의 열로 병기된다.

\subsection{정간보 전산화 연구}

정간보를 컴퓨터로 다루려는 시도는 크게 \emph{인코딩},
\emph{인식}, \emph{도구}의 세 갈래로 정리할 수 있다.

\textbf{인코딩.}
이용주~\cite{lee2013xml}는 정간보 기호를 분류하고 XML DTD로 악보 기술
방법을 정의한 뒤 이를 정간보 형태로 해석·표시하는 프로그램을 구현하여,
정간보의 구조적 기술이 가능함을 보였다. 이 방식은 후속 특허
~\cite{patent2017jgb}로 구체화되었다. XML 기술은 기계 처리에 적합하지만
여는·닫는 태그의 장황함 때문에 사람이 직접 읽고 쓰는 용도로는 부적합하다
(3.4절에서 정량 비교). 한편 Han 등~\cite{han2024sixdragons}은 세종실록악보
등 15세기 궁중음악을 대상으로, 분박 위치를 토큰 위치로 나타내는 정간보
충실형 인코딩을 제안하고 트랜스포머 언어모델로 선율 생성을 시연했다.
이는 정간보의 위치 기반 시가 표현이 기계학습 인코딩으로도 유효함을
보인 것으로, 본 연구의 텍스트 문법과 문제의식을 공유한다.

\textbf{인식.}
Kim 등~\cite{kim2025omr}은 정간보 인쇄 악보의 광학악보인식(OMR) 데이터셋과
인식 모델을 구축하여 기존 아날로그 악보의 대량 데이터화 경로를 열었다.
인식 결과를 담을 표준적·가독적 저장 포맷이 필요하다는 점에서
본 연구의 문법·교환 포맷과 상보적이다.

\textbf{도구.}
웹 기반 정간보 프로그램의 설계 연구~\cite{jkca2022web}와 공개 웹
편집기~\cite{jeongganbo-editor-web}가 있다. 후자는 격자에 율명을 넣고
PNG·JSON으로 내보내는 기본 기능을 제공하지만, 시김새·장단·가사 열,
대강 분절, 조판 세부 조정(칸 크기·간격·머리단), 인쇄 지면 관리, 재생 등
실제 악보 출판에 필요한 기능 범위에는 이르지 못했다.
\TODO{jkca2022web(한국콘텐츠학회논문지 2022, 웹 기반 정간보 프로그램 설계)의
저자·기능 범위를 원문으로 확인하고 비교 서술을 보강할 것}

\subsection{서양 기보 인코딩·웹 렌더링과의 비교}

서양 음악 정보 처리에서는 MusicXML~\cite{good2001musicxml}과
MEI~\cite{mei}가 교환·보존 포맷으로, Verovio~\cite{verovio}와
VexFlow 등이 웹 렌더링 엔진으로 정착했고, ABC~\cite{abc}나
Humdrum \code{**kern}~\cite{humdrum}처럼 사람이 타이핑하는 경량 텍스트
표기도 오래 쓰여 왔다. 본 연구의 문법은 이 가운데 ABC·\code{**kern}의
계보 --- ``타이핑 가능한 악보'' --- 를 정간보에 이식한 것으로 볼 수 있다.
다만 이들 포맷은 모두 음길이를 명시적 수치(예: \code{4분음표=4})로
기록하는 반면, 제안 문법은 정간보 원리 그대로 \emph{구분자와 공백의 위치}가
시가를 결정하도록 설계하여 악보 지면과 텍스트가 구조적으로 동형(isomorphic)이
되게 했다는 점이 다르다. 우종서 세로 조판 역시 기존 웹 렌더링 엔진이
지원하지 않아 자체 SVG 레이아웃 엔진을 구현했다(4.2절).
